\section{Conclusion}

Multi-tool selection is not a small variant of single-tool routing. Once one tool has been emitted, the model needs coverage, not repeated amplification of the same facet. The experiments in this paper show that a simple ontology-guided query-space bias has the correct sign on that problem, while a stationary key-side bias on the same basis does not.

The empirical picture is compact. On MetaTool Subtask4, Q-bias improves F1 on both Qwen and Llama, the ontology basis is necessary for the effect, and the stable K-side direction is not the smallest perturbation but the most specific one. These results establish a clear mechanism-level conclusion: in multi-tool selection, ontology-guided query suppression is aligned with coverage, whereas stationary key amplification is not.
